\documentclass[11pt,letterpaper]{article}
\usepackage[utf8]{inputenc}

%----- Configuración del estilo del documento------%
\usepackage{epsfig,graphicx}
\usepackage[left=2cm,right=2cm,top=1.8cm,bottom=2.3cm]{geometry}
\usepackage{fancyhdr}
\usepackage{lastpage}
\usepackage{url}
\pagestyle{fancy}
\fancyhf{}
\rfoot{\textit{Página \thepage \hspace{1pt} de \pageref{LastPage}}}


%------ Paquetes matemáticos básicos --------%
\usepackage{amsmath}
\usepackage{amssymb}
\usepackage{amsthm}

\usepackage[spanish]{babel}
\usepackage{graphicx}
\usepackage{hyperref}

\usepackage{tabularx}
\usepackage{xcolor}
\usepackage[table]{xcolor}
\usepackage{colortbl}
\usepackage{array, multirow, multicol, tabularx}
\usepackage{tcolorbox}
\newtheorem{theorem}{Theorem}[section]
\newtheorem{corollary}{Corollary}[theorem]
\newtheorem{lemma}[theorem]{Lemma}

%------si-------%
\definecolor{B}{HTML}{FFFFFF}
\definecolor{G}{HTML}{5e5e5e}
\definecolor{R2}{HTML}{d53d40}
\definecolor{A2}{HTML}{034190}
\definecolor{V2}{HTML}{7faa50}
\newcommand{\R}{\mathbb{R}}
\newcommand{\C}{\mathcal{C}}
\newcommand{\N}{\mathbb{N}}
\newcommand{\Z}{\mathbb{Z}}
\newcommand{\Q}{\mathbb{Q}}
\renewcommand{\theenumi}{\Roman{enumi}}
\renewcommand{\labelenumi}{{\theenumi}.}

\begin{document}

%------ Encabezado -------- %
\begin{center}
    \begin{minipage}{3cm}
    	\begin{center}
    		\includegraphics[height=3.4cm]{logo_unam.png}
    	\end{center}
    \end{minipage}\hfill
    \begin{minipage}{10cm}
    	\begin{center}
    	\textbf{\large Universidad Nacional Autónoma de México}\\[0.1cm]
        \textbf{Facultad de Ciencias}\\[0.1cm]
        \textbf{C\'alculo III}\\[0.1cm]
        Ejercicios lectura 1\\[0.1cm]
         El\'ias L\'opez Rivera\\[0.1cm]
        \texttt{ elias.lopezr\,@ciencias.unam.mx }\\[0.1cm]
        Fecha:\,\,28/07/2025
    	\end{center}
    \end{minipage}\hfill
    \begin{minipage}{3cm}
    	\begin{center}
    		\includegraphics[height=3.4cm]{Logo_FC.png}
    	\end{center}
    \end{minipage}
\end{center}

\rule{17cm}{0.1mm}

%------ Fin de encabezado -------- %
\,\\
\begin{tcolorbox}[
	title = \textcolor{black}{\textcolor{white}{Problema 2}}]
\textit{Sea $\chi$ un conjunto no vacio, decimos que una colecci\'on $\zeta\subset 2^{\chi} $ cumple la propiedad
de la intersecci\'on finita cuando para cada $A\subset \zeta$ de cardinalidad finita se cumple que:\,\\
\,\\
\begin{equation*}
    \bigcap_{B\in A}\,B \neq \emptyset
\end{equation*}\,\\
Demostrar que si $A\subset \R^n$ es compacto , entonces para cada $\chi\subset 2^A$ que cumpla la propiedad de la intersecci\'on finita
y que para cada $B\in \chi$ existe $U_{B}\in \R^n$ conjunto cerrado tal que $B=U_{B}\cap A$, se cumple que:\,\\
\,\\
\begin{equation*}
    \bigcap_{C\in\chi}\,C\neq \emptyset
\end{equation*}
}
\end{tcolorbox}
\begin{proof}\,\\
    \,\\
Procedemos por contradicci\'on, es decir existe $\chi\subset 2^A$ que cumple la propiedad de la intersecci\'on finita y que para todo $B\in \chi$ existe $U_B\in \R^n$ cerrado tal que $B=U_B\cap A$ que ademas cumple con:\,\\
\,\\
\begin{equation*}
    \bigcap_{B\in\chi}\,B=\emptyset
\end{equation*}\,\\
Si sacamos el complemento respecto a $A$:\,\\
\,\\
\begin{equation*}
    A=\emptyset^c=\left(\bigcap_{B\in\chi}\,B\right)^c=\bigcup_{B\in\chi}\,B^c=\bigcup_{B\in \chi}\,(U_B\cap A)^c
\end{equation*}\,\\
Como $U_B$ es cerrado en $\R^n$ y $A$ es cerrado en $\R^n$ se tiene que $U_B\cap A$ es cerrado en $\R^n$ y por tanto $(U_B\cap A)^c$ es abierto en $\R^n$ para todo $B\in \chi$, por tanto $A=\bigcup_{B\in \chi}\,(U_B\cap A)^c$ es una cubierta abierta de $A$, como 
$A$ es compacto existe $\zeta \subset \chi$ de cardinalidad finita tal que:\,\\
\,\\
\begin{equation*}
    A=\bigcup_{B\in \zeta}\,(U_B\cap A)^c
\end{equation*}\,\\
Sacando de nuevo el complemento respecto a $A$:\,\\
\,\\
\begin{equation*}
    \emptyset=A^c=\left(\bigcup_{B\in \zeta}\,(U_B\cap A)^c\right)^c=\bigcap_{B\in\zeta}\,(U_B\cap A)=\bigcap_{B\in \zeta}\,B
\end{equation*}\,\\
Pero como $\chi$ cumple la propiedad de la intersecci\'on finita y $\zeta\subset \chi$ de cardinalidad finita entonces:\,\\
\,\\
\begin{equation*}
    \emptyset\neq\bigcap_{B\in \zeta}\,B=\emptyset
\end{equation*}\,\\
Una contradicci\'on, como nuestra unica suposici\'on fue la existencia de $\chi$ concluimos la tesis del problema.
\end{proof}
\end{document}